% auto-generated by plot_kpi_full_test.py
\begin{table}[H]
\centering
\caption{全测试集（211 棵）学习策略剪枝前后 KPI 统计}
\label{tab:kpi_change_full_test}
\small
\setlength{\tabcolsep}{4pt}
\begin{tabular}{lcccc}
\toprule
\textbf{指标} & \textbf{剪前} & \textbf{剪后} & \textbf{变化量} & \textbf{改善占比 (\%)} \\
\midrule
结构层次性 $H$ & 0.4996 $\pm$ 0.0008 & 0.5000 $\pm$ 0.0009 & +0.0004 $\pm$ 0.0004 & 90.5 \\
拥挤度 $C$ & 0.4990 $\pm$ 0.0042 & 0.4988 $\pm$ 0.0041 & -0.0002 $\pm$ 0.0003 & 79.1 \\
平衡性 $B$ & 0.8342 $\pm$ 0.0963 & 0.8349 $\pm$ 0.0960 & +0.0006 $\pm$ 0.0005 & 96.7 \\
光照可及性 $\mathcal{L}$ & 0.3336 $\pm$ 0.0505 & 0.3351 $\pm$ 0.0512 & +0.0015 $\pm$ 0.0039 & 88.2 \\
果实潜力 $P$ & 0.4125 $\pm$ 0.0529 & 0.4157 $\pm$ 0.0530 & +0.0032 $\pm$ 0.0015 & 96.2 \\
结构稳定性 $S$ & 0.7018 $\pm$ 0.0382 & 0.7018 $\pm$ 0.0381 & -0.0000 $\pm$ 0.0011 & 85.8 \\
\bottomrule
\end{tabular}
\vspace{2pt}
\footnotesize 注：变化量为剪后减剪前；改善占比指朝期望方向变化的测试树比例；
$C$ 以减小为改善，其余指标以增大为改善。
\end{table}
